\chapter{Detailed Design}
\label{ch:detailed}

The detailed design decomposes the benchmark system into modules and describes
the internal behavior of each. This chapter presents the structure chart, then
the design of the validation modules, the negative builder, the mutator suite,
the metrics module, and the detector runners. It closes with a summary of the
design specifications.

%----------------------------------------------------------

\section{Structure Chart}

The system divides into three subsystems, each a group of cooperating modules.
The \textbf{corpus subsystem} holds the scraper, the sink-presence filter, the
audit driver, the diff filter, the negative builder, and the splitter. The
\textbf{mutation subsystem} holds the seven mutators and the
semantics-preservation checker. The \textbf{evaluation subsystem} holds the
detector runners, one per tool, and the metrics module. Control flows top down:
the corpus subsystem produces the validated split, which the mutation subsystem
consumes to produce variants, which the evaluation subsystem consumes alongside
the clean split to produce scores. Data flows back up only as artifacts on
disk, so no module holds state across runs.

\section{Module Description}

\subsection{Validation Modules}

The \textbf{sink-presence filter} defines, for each CWE, a small set of
high-recall regex sink patterns. It admits a sample only if at least one pattern
matches the source after comments and docstrings are stripped, which suppresses
documentation matches. The patterns match both safe and unsafe forms of the
sink, leaving discrimination to the diff filter. Two CWEs use proximity gates:
CWE-22 requires an HTTP-request reference near the sink.

The \textbf{audit driver} presents CWE-stratified samples to an independent
reviewer. For each item it supplies the identifier, source, repository, file
path, the matched regex, and the code excerpt around the sink, and it records a
binary correctly-labeled or mislabeled verdict with a one-sentence
justification. Reviewers do not see each other's verdicts, so the audit is
adversarial by construction.

The \textbf{diff filter} implements the evidence-anchored criterion. It compares
the vulnerable and fixed versions of a positive after comment stripping and
whitespace normalization, and keeps the positive only if at least one line
containing a sink-pattern match differs between the two versions. A positive
whose sink line is unchanged is rejected, because the fix did not touch the
sink and the label is therefore unsupported.

\subsection{Negative Builder}

The \textbf{negative builder} constructs the safe hard negatives. For each
negative it takes a real vulnerable function and applies the canonical
remediation for its CWE at the AST level, for example rewriting
\texttt{yaml.load} to \texttt{yaml.safe\_load} or a string-formatted SQL query
to a parameterized query. The remediation preserves the surface structure and
the sink \emph{name} while neutralizing only the sink's behavior. This gives
each negative a provable safe label and makes it a minimal-pair counterfactual:
a detector that fires on the mere presence of a sink token is forced into a
false positive on the twin.

\subsection{Mutator Suite}

The \textbf{mutator suite} generates seven semantics-preserving variants of each
test sample. Every mutator follows three design principles. It is
semantics-preserving: when a transformation cannot be proven safe, the mutator
skips that node. It is targeted at one shortcut, so a per-mutator accuracy drop
is diagnostic. And it is realistic: identifiers become synonyms and injected
code reads like ordinary scaffolding. The mutators fall into three families,
shown in Figure~\ref{fig:mutator-tax}.

\begin{figure}[htb]
\centering
\resizebox{0.95\columnwidth}{!}{%
\begin{tikzpicture}[
  font=\footnotesize,
  root/.style={draw, fill=gray!18, align=center, text width=15mm, minimum height=8mm},
  fam/.style={draw, rounded corners, fill=gray!10, align=center, text width=16mm, minimum height=6mm},
  leaf/.style={draw, align=left, text width=44mm, minimum height=5mm, inner sep=3pt},
  arr/.style={-{Stealth[length=1.8mm]}, gray!60},
]
\node[leaf] (m1) {M1\ \ dead-code injection};
\node[leaf, below=1mm of m1] (m2) {M2\ \ string-literal split$^{\dagger}$};
\node[leaf, below=1mm of m2] (m3) {M3\ \ identifier rename};
\node[leaf, below=1mm of m3] (m4) {M4\ \ wrapper extraction};
\node[leaf, below=5mm of m4] (m5) {M5\ \ sink attr obfuscation$^{\dagger}$};
\node[leaf, below=1mm of m5] (m6) {M6\ \ sink via builtins$^{\dagger}$};
\node[leaf, below=5mm of m6] (m7) {M7\ \ taint-through-dict$^{\dagger}$};
\node[fam, anchor=east] (f1) at ($(m1.west)!0.5!(m4.west) + (-9mm,0)$) {sink-\\preserving};
\node[fam, anchor=east] (f2) at ($(m5.west)!0.5!(m6.west) + (-9mm,0)$) {sink-\\targeted};
\node[fam, anchor=east] (f3) at ($(m7.west) + (-9mm,0)$) {dataflow};
\node[root, anchor=east] (r) at ($(f2.west) + (-10mm,0)$) {7\\mutators};
\draw[arr] (r.east) -- (f1.west);
\draw[arr] (r.east) -- (f2.west);
\draw[arr] (r.east) -- (f3.west);
\foreach \l in {m1,m2,m3,m4} \draw[arr] (f1.east) -- (\l.west);
\foreach \l in {m5,m6} \draw[arr] (f2.east) -- (\l.west);
\draw[arr] (f3.east) -- (m7.west);
\end{tikzpicture}}
\caption{The seven mutators grouped by the detector shortcut each family probes.
Mutators marked $\dagger$ are held out of training-time augmentation and used
only as test-time robustness probes.}
\label{fig:mutator-tax}
\end{figure}

The \emph{sink-preserving} family (M1 to M4: dead-code injection, string
splitting, identifier rename, wrapper extraction) alters code around the sink
but leaves the sink identifier and its arguments intact, probing positional,
substring, identifier-memorization, and single-function shortcuts. The
\emph{sink-targeted} family rewrites the sink call itself: M5 turns
\texttt{os.system(x)} into \texttt{getattr(os,"system")(x)} and M6 turns
\texttt{eval(x)} into \texttt{\_\_import\_\_("builtins").\_\_dict\_\_["eval"](x)},
so the sink token survives only inside a string literal. The \emph{dataflow}
family (M7) routes the sink's first argument through a one-entry dict,
\texttt{f(\{"a":x\}["a"])}, breaking the static lineage from parameter to sink
argument. At evaluation time two to three mutators are applied per sample, and
each mutator is also reported in isolation.

\subsection{Metrics Module}

The \textbf{metrics module} scores each detector one-versus-rest per CWE and
computes five bounded measures: macro-F1 across the seven classes, detection MCC
on the binary vulnerable-versus-safe collapse, Severity-Weighted Recall on the
CVSS-scored positives, Hierarchical CWE Macro-Accuracy using Wu-Palmer partial
credit on a hand-coded CWE subtree, and weighted Cohen's $\kappa$. It also
computes the robustness drop as clean macro-F1 minus mutated macro-F1 over the
same positives.

\subsection{Detector Runners}

Each \textbf{detector runner} adapts one tool to a common interface: it takes a
sample, invokes the tool, and returns a per-CWE verdict. The SAST runners invoke
Bandit and Semgrep as subprocesses and parse their findings. The whole-project
runners invoke Snyk Code and CodeQL. The LLM runners send a fixed prompt to each
model API at temperature zero. The trained runner loads the fine-tuned
GraphCodeBERT checkpoint and returns its predicted class. A uniform interface
lets the harness score every tier with the same metrics module.

\section{Summary}

The detailed design organizes the system into three subsystems of cooperating
modules: validation, mutation, and evaluation. The validation modules establish
label quality, the negative builder and mutator suite create the safe twins and
the robustness probes, and the metrics module and detector runners produce
comparable scores. Chapter~\ref{ch:impl} describes how these modules are
implemented.
